% CCSS 7.SP.C.7 — Experimental Probability (FL/TX: simulation emphasis)
\section{Experimental Probability}

\begin{learningGoals}
\begin{itemize}
    \item Approximate probability by collecting data from experiments
    \item Design and use simulations to estimate probabilities
\end{itemize}
\end{learningGoals}

\begin{conceptBox}{Experimental Probability and Simulations}
\textbf{Experimental probability} uses real data:
$$P(\text{event}) = \frac{\text{number of times event occurs}}{\text{total trials}}$$
A \textbf{simulation} models a real situation using simple tools like coins, dice, spinners, or random-number generators. Simulations let you estimate probabilities when real experiments are too slow, expensive, or impractical.

\textbf{Steps to design a simulation:}
\begin{enumerate}[leftmargin=*]
    \item Identify the event and its possible outcomes.
    \item Choose a model (coin, die, random digits) that matches the probabilities.
    \item Run many trials and record the results in a frequency table.
    \item Calculate the experimental probability and compare to the theoretical value.
\end{enumerate}
\end{conceptBox}

\begin{workedExample}{Simulating Blood Type Donors}
About $30\%$ of donors have type~A blood. Use random digits ($0$--$9$) to simulate finding the first type~A donor. Let digits $0, 1, 2$ represent type~A.

\begin{center}
\begin{tabular}{c c c}
\textbf{Trial} & \textbf{Random digits} & \textbf{Donors until type A} \\
\hline
$1$ & $7, 4, \mathbf{1}$ & $3$ \\
$2$ & $\mathbf{0}$ & $1$ \\
$3$ & $8, 5, 6, \mathbf{2}$ & $4$ \\
$4$ & $3, \mathbf{1}$ & $2$ \\
$5$ & $9, 7, \mathbf{0}$ & $3$ \\
\end{tabular}
\end{center}
Average donors needed $= (3 + 1 + 4 + 2 + 3) \div 5 = 2.6$.
\answerHighlight{The simulation suggests about $2$--$3$ donors are needed to find one with type~A blood.}
\end{workedExample}

\mascotSays{More trials = better estimates. $10$ trials give a rough idea; $100$+ trials give reliable results.}

\begin{practiceBox}{Experimental Probability}
\resetProblems

\prob A coin is tossed $100$ times. Heads comes up $48$ times. What is the experimental probability of heads? How does it compare to the theoretical probability?
\answerExplain{$0.48$; close to the theoretical $0.50$}{$P = 48 \div 100 = 0.48$. The theoretical probability is $0.50$. The small difference is expected — with more tosses, the experimental value would get closer to $0.50$.}

\prob A spinner has $4$ equal sections: red, blue, green, yellow. After $80$ spins, red appeared $24$ times. Find the experimental probability. Is the spinner fair?
\answerExplain{$0.30$; slightly above $0.25$ but likely fair}{$P = 24 \div 80 = 0.30$. Theoretical $= 0.25$. The small difference could be due to chance in $80$ trials. More spins would help confirm.}

\prob Design a simulation to estimate the probability that a family with $3$ children has all girls. Describe your model and how to run it.
\answerExplain{Use a coin; H = boy, T = girl; flip $3$ times per trial}{Each child is equally likely to be a boy or girl — like a coin flip. Flip $3$ coins per trial. Record how many trials give $3$ tails (all girls). Repeat $50$+ trials. Divide all-girl trials by total trials.}

\prob In a simulation, a die is rolled $60$ times. A $6$ appeared $15$ times. What is the experimental probability? What would you expect theoretically?
\answerExplain{$0.25$; theoretical is $\frac{1}{6} \approx 0.167$}{$P = 15 \div 60 = 0.25$. Theoretical $= \frac{1}{6} \approx 0.167$. The experimental value is noticeably higher — this could be chance, or the die may not be fair.}

\prob A random-number generator produces digits $0$--$9$. You want to simulate a $40\%$ chance of rain. Which digits would represent rain? Run $10$ trials with: $3, 7, 1, 0, 9, 4, 2, 8, 5, 6$. What experimental probability do you get?
\answerExplain{Rain = digits $0$--$3$; experimental $P = 0.40$}{$40\% = 4$ out of $10$ digits. Choose $0, 1, 2, 3$ for rain. In the $10$ trials: rain digits are $3, 1, 0, 2$ — that's $4$ rain events. $P = 4 \div 10 = 0.40$.}

\end{practiceBox}
